\documentclass[12pt,a4paper]{article}
\usepackage[utf8]{inputenc}
\usepackage[bahasa]{babel}
\usepackage[margin=3cm]{geometry}
\usepackage{setspace}
\usepackage{parskip}
\usepackage{fancyhdr}

\onehalfspacing

\begin{document}

\begin{flushright}
Makassar, 19 September 2026
\end{flushright}

Hal \> : \> Gugatan Perbuatan Melawan Hukum (PMH)\\
Lampiran \> : \> Surat Kuasa\\


\vspace{0.5cm}

Kepada Yth.\\
\textbf{Ketua Pengadilan Negeri Makassar}\\
di - \\
\quad \textbf{Makassar}

Dengan hormat,
Yang bertanda tangan dibawah ini:
\begin{enumerate}
    \item Liya
    \item Liya
    \item Liya
\end{enumerate}
Masing masing adalah Advokat dari Kantor "Firma Hukum Berkeadilan" Jalan Beringin, No.7.Berdasarkan Surat Kuasa Khusus Nomor: 198/SKK-PDTG/FHB/IX-2026 pada tanggal 15 September 2026 (terlampir), bertindak atas nama \textbf{Rahmat Hidayat} yang berkedudukan hukum di Jalan Sukaribut, Kota Makassar.
Selanjutnya disebut \textbf{PENGGUGAT}

Dengan ini, mengajukan gugatan perbuatan melawan hukum dan tuntutan ganti rugi terhadap:
\begin{tabbing}
    Nama \quad\quad\quad\quad \= : \textbf{Joko Saputra}\\
    Alamat \> : Jalan Sukaribut, Kota Makassar\\
\end{tabbing}
Selanjutnya disebut sebagai \textbf{TERGUGAT};

Gugatan ini diajukan atas alasan-alasan sebagai berikut:

\begin{enumerate}
    \item Bahwa Penggugat adalah pemilik sah dan telah menempati sebuah rumah sederhana di Kota Makassar selama lebih dari 8 (delapan) tahun. Di bagian samping rumah Penggugat terdapat jendela dan ventilasi udara yang menghadap ke tanah kosong.
    \item Bahwa ventilasi tersebut selama bertahun-tahun berfungsi sebagai sumber utama sirkulasi udara dan cahaya alami bagi ruang keluarga dan kamar tidur rumah Penggugat.
    \item Bahwa pada tahun 2024, tanah kosong di sebelah rumah Penggugat dibeli oleh Tergugat, yang kemudian mendirikan bangunan rumah dua lantai di atas tanah tersebut.
    \item Bahwa Tergugat membangun tembok setinggi kurang lebih 4 (empat) meter tepat di batas tanah yang sangat dekat dengan dinding rumah Penggugat, sehingga menutup total ventilasi dan jendela rumah Penggugat.
    \item Bahwa akibat tertutupnya ventilasi tersebut, rumah Penggugat menjadi gelap, pengap, lembap, serta berdampak buruk pada kesehatan anak Penggugat yang memiliki riwayat penyakit pernapasan.
    \item Bahwa tindakan Tergugat yang membangun tembok tanpa mengindahkan hak-hak tetangga, menolak beriktikad baik, serta menolak bermusyawarah merupakan bentuk \textbf{Perbuatan Melawan Hukum (Onrechtmatige Daad)} sebagaimana diatur dalam Pasal 1365 KUHPerdata.
\end{enumerate}

Berdasarkan alasan-alasan tersebut di atas, Kami mohon Ketua Pengadilan Negeri Makassar yang memeriksa dan mengadili perkara ini berkenan memberikan putusan sebagai berikut:

\textbf{PRIMER:}
\begin{enumerate}
    \item Mengabulkan gugatan Penggugat untuk seluruhnya;
    \item Menyatakan demi hukum bahwa tindakan Tergugat yang membangun tembok sehingga menutup total ventilasi dan jendela rumah Penggugat adalah \textbf{Perbuatan Melawan Hukum};
    \item Memerintahkan Tergugat untuk membongkar atau menurunkan sebagian tembok yang menutup ventilasi udara dan jendela rumah Penggugat selambat-lambatnya 7 (tujuh) hari sejak putusan ini berkekuatan hukum tetap;
    \item Menghukum Tergugat untuk membayar ganti rugi materiel dan immateriel kepada Penggugat;
    \item Menghukum Tergugat untuk membayar segala biaya yang timbul dalam perkara ini.
\end{enumerate}

\textbf{SUBSIDER:}

Mohon putusan yang seadil-adilnya (\emph{ex aequo et bono}).

Demikian gugatan ini Kami ajukan, atas perhatian dan dikabulkannya gugatan ini diucapkan terima kasih.


\vspace{1cm}

\begin{flushright}
    Hormat kami,\\
    Kuasa Hukum Penggugat,\\[1.5cm]
    Liya
\end{flushright}

\end{document}
